\printedpage{4}
% head: A. Grothendieck : Dérivateurs, Ch. XII, Édition des Universités de Montpellier II et Paris VII
We use the proof of I, p. 75, with the help of the category called $\Phi(X)$ ($\mathrm{Ob}\,\Phi(X) = \mathrm{Fl}(X)$, but with a somewhat peculiar notion of arrow in $\Phi(X)$). We find that $f \in W \iff \Phi(f) \in W$, and likewise $f^\circ \in W \iff \Phi(f) \in W$.

NB. In the proof of proposition 2, as in that of proposition 1 (for the implication b) $\Longrightarrow$ a)), we used the localization property W(2) only in the weak form of the corollary, i.e. in the form $f$ $W$-aspheric $\Longrightarrow f \in W$.

\textbf{Corollary} (of prop. 2): \textit{Let $W \subseteq \mathrm{Fl}(\mathrm{Cat})$ satisfy the axiom} W(1) (saturation). \textit{Then the conjunction of the axioms} W(2), W(3) \textit{is equivalent to that of the dual axioms} W($2'$), W($3'$) (\footnote[5]{Likewise if in this statement we replace W(2) by W(2 bis).}).

It suffices to prove that 1,2,3 $\Longrightarrow$ $2'$,$3'$.

\msp{page 4}

For $2'$ \texttt{[rather} $3'$\texttt{]}, this is corollary 2 above. For $3'$ \texttt{[rather} $2'$\texttt{]}, this follows at once from proposition 2. (NB This statement would remain valid, if we would have taken W(2) in the weak form only, cf. the NB on the previous page.)

I still have to put in an axiom of the form L5, L6 (cf. VI, p. 43-44). But the statement of these axioms uses the notion of Hot-fibered objects, or of perfect functors, i.e. calls on knowledge of $W_\infty$, which I would like to avoid (if I can). I must therefore introduce the notion of $W$-fibration, on the model of what I did in the context of derivators, in I (cf. p. 111-126). We introduce the notion of \textbf{D}-fibration on page 117 (prop. 5), by \texttt{[some]} equivalent conditions, some of which make sense in terms of the datum of $W$ alone. So let
\[ f : X \longrightarrow Y \]
be in Cat and consider the conditions

\begin{enumerate}
\item[a)] If $Y'' \longrightarrow Y'$ in $\mathrm{Cat}/Y$ is in $W$, then so is $X'' \longrightarrow X'$ (base change by $X \longrightarrow Y$ takes $W$ into $W$).
\item[a$'$)] As in a), but with $Y''$, $Y'$ $W$-aspheric.
\item[a$_1$)] As in a), in the particular case where $Y'$ has a final object and $Y''$ is the one-point category.
\item[a$_2$)] Dual of a$_1$), with an initial object of $Y'$.
\end{enumerate}

We thus have
\[
\begin{tikzcd}[row sep=large, column sep=1.4em]
 & & \mathrm{a} \arrow[d, Rightarrow] & \text{ultrastrong $W$-fibration} & \\
 & & \mathrm{a}' \arrow[dl, Rightarrow] \arrow[dr, Rightarrow] & \text{strong $W$-fibration} & \\
\text{left $W$-fibration} & \mathrm{a}_1 & & \mathrm{a}_2 & \text{right $W$-fibration}
\end{tikzcd}
\]
% foot: Version 7 décembre 2002
\printedpage{5}
% head: A. Grothendieck : Dérivateurs, Ch. XII, Édition des Universités de Montpellier II et Paris VII
and in the case $W = W_\infty$ all these conditions are equivalent. I have a mind \texttt{[to]} take the strongest of these conditions, but I must give myself some means of finding $W$-fibrations in so strong a sense. For example I would like to know that if $X$
\msp{page 5}
is in Cat, then $\Ch_\infty(X) \longrightarrow X \times X$ or $\Ch(X) \longrightarrow X \times X$ is a $W$-fibration. In both cases, and (say) for $W_\infty$, I prove it by proving that $p$ is both smooth, and proper. So we would have to go over the notions of smoothness and properness again in the context of the datum of a $W$.

If I remember correctly, property a$'$) \texttt{[rather a}$_1$\texttt{)]} was something fairly easily checked, it sufficed to suppose that the
\[ X_y \longrightarrow X/y \qquad \text{and} \qquad X/y \longrightarrow X/y' \]
(for $y \in \mathrm{Ob}\,Y$ resp. $u : y \longrightarrow y'$ in $\mathrm{Fl}\,Y$) are in $W$, and that this remains true under every base change - which amounts to condition a$_1$) on $f$. We could therefore formulate a condition on $W$ as follows:

\begin{enumerate}[leftmargin=*]
\item[\textbf{W(4)}] (Axiom of $W$-fibrations (\footnote[6]{See sufficient weakened variant, in the NB p. 10. See also [?] of variants pages 13, 14 below.}).) Let $f : X \longrightarrow Y$ in Cat satisfy the $W$-fibration condition a$_1$) above (p. 4), then it even satisfies a), i.e. base change by $f$ takes arrows in $W$ to arrows in $W$ (\footnote[7]{i.e. a left $W$-fibration is an ultrastrong fibration.}).
\end{enumerate}

Thanks to proposition 2, it follows that the dual condition a$_2$) on $f$ implies the strong $W$-fibration condition, so the axiom W(4), apparently asymmetric, also means that the four fibration conditions on $f$ are equivalent, it is therefore in fact symmetric.

\textbf{NB} I do not see, without additional axioms on $W$, how to pass from a$_1$ to a$'$ alone. The difficulty is to show, by means of a$_1$ alone, that if $Y$ is $W$-aspheric, the inclusion $X_y \longrightarrow X$ (where $y \in \mathrm{Ob}\,Y$) is in $W$. In the context of derivators \textbf{D}, this could be seen by \textbf{D}-cohomological means.

This very strong axiom suffices to imply that if $X$, $Y$ are $W$-fibered over $S$, and if $f : X \longrightarrow Y$ is an $S$-morphism which is in $W$, then $f_s : X_s \longrightarrow Y_s$ is in $W$ ($\forall\ s \in S$). Let us at least show the converse for now.

\msp{page 6}

\textbf{Proposition 3.} \textit{Let $X$, $Y$ be in} Cat$/S$\textit{, $f : X \longrightarrow Y$ an $S$-arrow. Suppose ($X$, $Y$ $W$-fibered) that $\forall\ s \in S$, $f_s : X_s \longrightarrow Y_s$ is in $W$, then $f \in W$} (\footnote[8]{\textit{Without} axiom W(4). This is the first time we need W(2) in its strong form, to the exclusion of W(2 bis).}).

We have to use the localization axiom, hence look at the $f/s$ induced by $f$
% foot: Version 7 décembre 2002
\printedpage{6}
% head: A. Grothendieck : Dérivateurs, Ch. XII, Édition des Universités de Montpellier II et Paris VII
\[
\begin{tikzcd}[row sep=large]
X/s \arrow[r, "f/s"'] & Y/s \\
X_s \arrow[u, "i_s^X"] \arrow[r, "f_s"] & Y_s \arrow[u, "i_s^Y"'] \quad .
\end{tikzcd}
\]

Condition a$_1$ suffices to imply that $i_s^X$ and $i_s^Y$ are in $W$, so $f/s$ is in $W$ if and only if $f_s$ is. We did not have to use axiom W(4), and it sufficed to suppose that $X$, $Y$ are ``$W$-cohomological'' fibered objects (or left $W$-fibered) over $S$, or dually, right $W$-fibered objects.

Moreover:

\textbf{Corollary} (\footnote[9]{Without axiom W(4).}). \textit{Let $S_0 \subseteq \mathrm{Ob}\,S$ be a subset of $S$ such that $S_0 \longrightarrow \pi_0(S)$ surjective. Then in proposition 3 it suffices to suppose $f_s \in W$ for $s \in S_0$.}

It suffices to see that if $u : s \longrightarrow s'$ is an arrow in $S$, then $f_s \in W \iff f_{s'} \in W$. We use the fact that $f_s \in W \iff f/s \in W$, and likewise for $f_{s'}$. So to see that $f/s \in W \iff f/s' \in W$. Now we have
\[
\begin{tikzcd}[row sep=large]
X/s \arrow[r, "f/s"] \arrow[d, "X/u"'] & Y/s \arrow[d, "Y/u"] \\
X/s' \arrow[r, "f/s'"] & Y/s' \quad ,
\end{tikzcd}
\]
where $X/u$ and $Y/u$ are in $W$, whence the result.

But let us return to the converse: if $f \in W$, then do we have $f_s \in W_S$? \texttt{[rather} $f_s \in W$\texttt{?]} There is an argument to this effect in VI, p. 23, which I am going to take up again. The idea is this. We consider the connected component $Y_0$ of $Y$ \texttt{[rather} $S_0$ \texttt{of} $S$\texttt{]} containing $s$, and we construct a
\msp{page 7}
$W$-fibered object (in the \textit{strong sense} a))
\[ S' \longrightarrow S_0 \subseteq S \]
with $S'$ \textit{$W$-aspheric.} Since $X \longrightarrow Y$ over $S$ is in $W$, the same will therefore hold for $X' \longrightarrow Y'$ over $S'$. Since the fibers of $X'$, $Y'$ at $s'$ are identified with those of $X$, $Y$ at the image of $s'$ in $S$, we are reduced to the case where $S$ is $W$-aspheric. But then we have
\[
\begin{tikzcd}[row sep=large]
X \arrow[r, "f"] & Y \\
X_s \arrow[u, "i_s^X"] \arrow[r, "f_s"] & Y_s \arrow[u, "i_s^Y"'] \quad ,
\end{tikzcd}
\]
% foot: Version 7 décembre 2002
\printedpage{7}
% head: A. Grothendieck : Dérivateurs, Ch. XII, Édition des Universités de Montpellier II et Paris VII
where $i_s^X$, $i_s^Y$ are in $W$, since obtained from $\{s\} \longrightarrow S$ (which is in $W$, $S$ being $W$-aspheric) by the $W$-fibrant base changes $X \longrightarrow S$, $Y \longrightarrow S$. Hence $f \in W \implies f_s \in W$.

NB We used the fibration property for $X$, $Y$ only in the sense a$'$) (\footnote[10]{And even, in a$'$) we could require that $Y'' = e$, $Y'$ $W_\omega$-aspheric, where $W_\omega$ \texttt{[is]} defined on page 14, $W_\omega \subseteq W \cap W_\infty \ldots$}).

We must therefore prove, if we can, the

\textbf{Lemma} (?) (\footnote[11]{Is proved by means of axiom W(4) (+ some small later verifications).}). \textit{Let $S$ be a} $0$\textit{-connected object of} Cat\textit{. Then there exists a $W$-aspheric object $S'$ of} Cat\textit{, and a $W$-fibration (in the} strong \textit{sense of a), page 4) $S' \longrightarrow S$.}

The natural idea is to take, for fixed $s_0 \in \mathrm{Ob}\,S$,
\[ S' = \Ch(S; s_0, -) \qquad \text{or} \qquad S' = \Ch_\infty(S; s_0, -), \]
and the arrow
\[ \beta : c \longmapsto \text{target } c : S' \longrightarrow S. \]
We know that $S'$ is $W_\infty$-apheric and that $\beta$ is $W_\infty$-fibrant (and even perfect). We must see that $S'$ is also $W$-aspheric, and $\beta$ $W$-fibrant. We must therefore go over the proof again in the case of $W_\infty$, and prove that it transforms to $W$, % [sic] aphérique -> asphérique, lettre s omise
\msp{page 8}
given the axioms needed on $W$. I have a mind to look first at the choice $S' = \Ch(S; s_0, -)$, without passing to a direct limit which would force me to put a limit axiom on $W$. The delicate and essential homotopy properties of $\Ch(S)$ are established in VII, p. 157-169, namely that in the diagram
\[ S \xrightarrow{i_S} \Ch(S) \xrightarrow{p_S} S \times S, \]
$i_S \in W_\infty$ and $p_S$ is a $W_\infty$-fibration (in fact, \textit{proper} and \textit{smooth} for $W_\infty$) (\footnote[12]{In the statement of the ``essential properties'' of $\Ch(S)$, we can work with $W_\omega$, the smallest subset of Fl(Cat) satisfying axioms W(1), W(2 bis), W(3) (cf. page 14), so $W_\omega \subseteq W \cap W_\infty$.}). The fact that $i_S \in W$ \texttt{[rather} $i_S \in W_\infty$\texttt{]} implies that the source and target functors
\[ \begin{tikzcd} \Ch(S) \arrow[r, shift left, "\sigma"] \arrow[r, shift right, "\beta"'] & S \end{tikzcd} \qquad \sigma = \mathrm{pr}_1 \circ p_S,\ \beta = \mathrm{pr}_2 \circ p_S \]
are in $W_\infty$ (since $\sigma i_S = \beta i_S = \mathrm{id}_S$). Since these are Hot-fibrations, it follows that the fibers $\Ch(S; s, -)$ resp. $\Ch(S; -, s)$ are Hot-aspheric. But even granting that $\sigma$, $\beta$ are in $W$ and are $W$-fibrations, I cannot conclude right away that the fibers are $W$-aspheric - there would be a vicious circle! So I need a direct proof.

Consider the functor
\[
\begin{tikzcd}[row sep=huge]
\Ch(S; s, -) \arrow[d, "p"'] \\
\underline{Ty} \quad ,
\end{tikzcd}
\]
% foot: Version 7 décembre 2002
\printedpage{8}
% head: A. Grothendieck : Dérivateurs, Ch. XII, Édition des Universités de Montpellier II et Paris VII
it is a Cat-fibrant functor, hence $W$-smooth (I will have to develop the notion of $W$-smoothness (\footnote[13]{$f : X \longrightarrow Y$ is $W$-smooth if the $X_y \longrightarrow y\backslash X$ are $W$-aspheric.}) carefully, I will assume for the moment the familiar formal properties, and will justify them later). Because of $W$-smoothness, to prove $p \in W$, it suffices to
\msp{page 9}
prove that the fibers of $p$ are $W$-aspheric (\footnote[14]{Here we use W(2 bis), not $W(2)$.}). Now the fibers are the $Z_\tau = \Ch_\tau(S; s, -)$, which are \textit{contractible} categories (i.e. $Z_\tau \longrightarrow e$ is a homotopism), hence $W$-aspheric by W(3) (cf. cor. 3, page 3).

On the other hand, if we grant that $\Ch(S) \xrightarrow{p_S} S{\times}S$ is $W$-smooth (and $W$-proper), hence $\Ch(S) \xrightarrow{\sigma_S} S$ also, having $W$-aspheric fibers as we have just seen, it follows that $\sigma_S \in W$, whence $S \xrightarrow{i_S} \Ch(S)$ also in $W$ (since $\sigma_S \circ i_S = \mathrm{id}_S$). So it remains only to prove that
\[ \Ch(S) \xrightarrow{p_S} S \times S \tag*{$(*)$} \]
is $W$-proper and $W$-smooth. It will follow that
\[ \Ch(S; s_0, -) \xrightarrow{\beta} S \]
is likewise $W$-proper and $W$-smooth, since obtained from $p_S$ by the base change
\[ S \longrightarrow S \times S \qquad s \longmapsto (s_0, s). \]

Next we must see that $W$-smooth and $W$-proper implies the $W$-fibration property \textit{in its strongest form} a) (whereas I foresee that a priori we will only have it in the weak forms a$_1$, a$_2$ - not even in the form a$'$). So it is here that the axiom W(4) of $W$-fibrations will be essential (\footnote[15]{in the weakened form W(4a$''$) (p. 13).})!

Here is the argument for the smoothness of $(*)$, in loc. cit. p. 161 ff. We must prove that for $c \in \mathrm{Ob}\,\Ch(S)$, and an arrow $u \in \mathrm{Fl}(S \times S)$ \texttt{[rather} $(u,v) \in$ \texttt{Fl(}$S \times S$\texttt{)]}, a certain category of ``liftings of $u$ with source $c$ '' \texttt{[rather ``liftings of} $(u,v)$ \texttt{with target} $c$ \texttt{'']}, $\mathcal{C} = \mathcal{C}(c,u,v)$, is $W$-aspheric (\footnote[16]{In fact, we prove that $\mathcal{C}(c,u,v)$ is \textit{contractible}, hence it is $W_\omega$-aspheric.}). For this, we compare it to a category $\mathcal{C}'$ with a final object, hence contractible and a fortiori $W$-aspheric, with the help of two functors
\[ \begin{tikzcd} \mathcal{C} \arrow[r, shift left, "\varphi"] & \mathcal{C}', \arrow[l, shift left, "\psi"] \end{tikzcd} \]
and the homotopies $\varphi\psi \sim \mathrm{id}_{\mathcal{C}'}$ (no merit, $\mathcal{C}'$ being contractible)
\msp{page 10}
and $\psi\varphi \sim \mathrm{Id}_{\mathcal{C}}$ (we even find an arrow $\psi\varphi \longrightarrow \mathrm{id}_{\mathcal{C}}$). Thus $\varphi$ is a homotopism, hence a $W$-equivalence (cor. 3, p. 3), hence $\mathcal{C}$ is $W$-aspheric, $\mathcal{C}'$ being so. The properness of $p_S$ is equivalent to the smoothness of $p_{(S^\circ)}$, it is therefore proved at the same time.
% foot: Version 7 décembre 2002